\chapter{Testing maya applications}
\label{ch:testing}

\begin{aside}
A library architecture is only worth the cost it imposes if the
benefits show up in tests. \maya{}'s purity discipline pays off
here: most of an app is testable without a terminal, a renderer,
or even a runtime. This chapter is a playbook of testing
patterns that work with the architecture.
\end{aside}

\section{What is and is not testable directly}

\code{update} is a pure function. So is \code{view}. So is the
fingerprint computation, the layout pass, and the diff. These
are testable with plain function calls and assertions; no setup.

The things that are \emph{not} directly testable are anything
that touches the OS:

\begin{itemize}[leftmargin=*]
  \item The event loop blocking on epoll.
  \item The renderer writing bytes to a terminal.
  \item Tasks spawning worker threads.
  \item Signals delivered to signalfd.
\end{itemize}

For these, you mock the boundary. \maya{} ships a test harness
that wraps the runtime with a controllable event source and
captured-byte renderer.

\section{Unit testing update}

The simplest test pattern:

\begin{lstlisting}
TEST_CASE(\"counter increments\") {
    auto m0 = Counter::init().first;
    auto [m1, cmd] = Counter::update(m0, Counter::Inc{});
    REQUIRE(m1.count == 1);
    REQUIRE(std::holds_alternative<None>(cmd.v));
}
\end{lstlisting}

Direct function call. No runtime, no terminal, no setup. The
test is the specification.

\subsection*{Testing a sequence}

A useful helper:

\begin{lstlisting}
template <Program P>
auto fold(std::initializer_list<typename P::Msg> msgs) {
    auto m = P::init().first;
    std::vector<Cmd<typename P::Msg>> cmds;
    for (auto& msg : msgs) {
        auto [m2, c] = P::update(m, msg);
        m = std::move(m2);
        cmds.push_back(c);
    }
    return std::pair{std::move(m), std::move(cmds)};
}

TEST_CASE(\"three increments and a decrement\") {
    auto [m, _] = fold<Counter>({
        Counter::Inc{}, Counter::Inc{}, Counter::Inc{},
        Counter::Dec{}
    });
    REQUIRE(m.count == 2);
}
\end{lstlisting}

The fold is the playback. The final model is the assertion target.

\subsection*{Testing emitted commands}

When \code{update} returns a non-trivial \code{Cmd}, assert
against the variant:

\begin{lstlisting}
TEST_CASE(\"save schedules a clipboard write\") {
    auto [m, cmd] = MyApp::update(
        MyApp::Model{.body = \"hello\"}, MyApp::Save{});
    auto* wc = std::get_if<WriteClipboard<MyApp::Msg>>(&cmd.v);
    REQUIRE(wc != nullptr);
    REQUIRE(wc->text == \"hello\");
}
\end{lstlisting}

\code{std::get\_if} returns a pointer or null; testing for null
is the cleanest pattern.

\section{Testing view}

\code{view} returns an \code{Element}. You can assert against
its structure with a visitor:

\begin{lstlisting}
TEST_CASE(\"counter view contains the count\") {
    auto e = Counter::view(Counter::Model{.count = 42});
    bool found = false;
    walk(e, [&](const TextElement& t) {
        auto s = text_view(t);
        if (s.find(\"42\") != std::string_view::npos) found = true;
    });
    REQUIRE(found);
}
\end{lstlisting}

\code{walk} is a tree-walker helper. You assert on the presence
or absence of text, not on the exact rendering. The exact bytes
depend on the terminal capabilities.

\subsection*{Snapshot tests}

For UIs that are mostly stable, snapshot tests work well:

\begin{lstlisting}
TEST_CASE(\"counter view snapshot\") {
    auto e = Counter::view(Counter::Model{.count = 42});
    auto dump = dump_element(e);
    REQUIRE(dump == load_snapshot(\"counter_view_42.txt\"));
}
\end{lstlisting}

\code{dump\_element} produces a textual representation; the
snapshot file is checked into source control. Differences are
visible in code review.

The risk: snapshot tests catch every change, including
intentional ones. Update them deliberately, not reflexively.

\section{Property tests}

Pure functions are gold for property-based testing:

\begin{lstlisting}
PROPERTY(\"increment then decrement is identity\", Counter::Model m) {
    auto [m1, _] = Counter::update(m, Counter::Inc{});
    auto [m2, __] = Counter::update(m1, Counter::Dec{});
    REQUIRE(m2.count == m.count);
}

PROPERTY(\"any sequence preserves count >= 0\", std::vector<Counter::Msg> seq) {
    auto m = Counter::init().first;
    for (auto& msg : seq) {
        auto [m2, _] = Counter::update(m, msg);
        m = std::move(m2);
    }
    REQUIRE(m.count >= 0);
}
\end{lstlisting}

The framework (rapidcheck, Catch2 generators) supplies the
inputs. Counterexamples are minimised automatically. Bugs that
would never appear in hand-written tests show up.

\subsection*{Generating Msg variants}

For a \code{Msg = std::variant<\ldots>}, the generator needs to
sample each alternative. Custom generators are straightforward:

\begin{lstlisting}
template <>
struct Arbitrary<Counter::Msg> {
    static Gen<Counter::Msg> arbitrary() {
        return gen::oneOf<Counter::Msg>(
            gen::just(Counter::Inc{}),
            gen::just(Counter::Dec{}),
            gen::just(Counter::Quit{})
        );
    }
};
\end{lstlisting}

For variants with payload, generate the payload with
\code{gen::build}.

\section{Testing the layout phase}

The layout pass is a pure function from
\code{(Element, Size)} to a layout tree:

\begin{lstlisting}
TEST_CASE(\"flex row distributes width\") {
    auto e = h(
        text(\"abc\") | grow<1>,
        text(\"def\") | grow<2>
    );
    auto layout = compute_layout(e, Size{30, 1});
    REQUIRE(rect_of_child(layout, 0).w == 10);
    REQUIRE(rect_of_child(layout, 1).w == 20);
}
\end{lstlisting}

The layout helper exposes the post-layout rects so you can
assert against them.

\section{Testing the diff and serialise stages}

These are pure too, but assertions become byte-counting:

\begin{lstlisting}
TEST_CASE(\"changing one cell emits a minimal diff\") {
    Canvas a = canvas_of(\"hello world\");
    Canvas b = canvas_of(\"hello WORLD\");
    auto diff = compute_diff(a, b);
    auto bytes = serialise(diff);
    REQUIRE(bytes.size() < 30);  // cursor move + 5 cells + reset
}
\end{lstlisting}

The exact byte count depends on the serialiser strategy; the
contract is ``less than full redraw''.

\section{End-to-end with a fake terminal}

For tests that exercise the runtime end-to-end:

\begin{lstlisting}
TEST_CASE(\"app handles quit message\") {
    test::FakeTerminal term;
    test::EventQueue events;
    events.push(Event::KeyPress{'q'});

    auto rc = test::run<Counter>(term, events);
    REQUIRE(rc == 0);
    REQUIRE(term.last_frame_contains(\"Count: 0\"));
}
\end{lstlisting}

The fake terminal records every byte written. The event queue
delivers events in order. The test runner exits when the program
returns. The terminal's \code{last\_frame\_contains} is a
helper for fuzzy assertions.

\subsection*{Time control}

Real-time-dependent tests use a virtual clock:

\begin{lstlisting}
test::FakeClock clock;
events.push_at(clock.now() + 1s, Event::KeyPress{'a'});

auto rc = test::run<MyApp>(term, events, clock);
\end{lstlisting}

The virtual clock advances only when the event queue asks for
a wait. Tests that need to verify ``300ms later, the toast
expired'' run instantly.

\section{Testing animation}

Animation tests use \code{AnimationFrame} subscriptions and
the virtual clock:

\begin{lstlisting}
TEST_CASE(\"spinner advances each frame\") {
    test::FakeTerminal term;
    test::EventQueue events;
    test::FakeClock clock;

    // Advance 10 frames.
    for (int i = 0; i < 10; ++i) {
        events.push_animation_frame();
    }

    test::run<App>(term, events, clock, /*max_frames=*/10);
    REQUIRE(term.frame_count() == 10);
    // Each frame should show a different spinner glyph.
}
\end{lstlisting}

The \code{max\_frames} parameter prevents infinite loops in
tests that would otherwise run forever.

\section{Testing concurrency}

\code{Cmd::task} runs on a thread. To test deterministically,
use a synchronous task pool:

\begin{lstlisting}
auto ctx = test::Context{.task_pool = test::SyncTaskPool{}};
\end{lstlisting}

Synchronous tasks run inline on the main thread; the assertion
about their result happens after the dispatch.

\subsection*{Testing race conditions}

Race-condition tests are the hardest. \maya{} ships a
\code{ThreadSanitizer}-aware build mode that flags races at
runtime; pair it with the test suite by running tests with
\code{-fsanitize=thread}.

\begin{warning}
Race tests are necessarily probabilistic. A test that passes
once does not prove the absence of races; it only proves the
race did not happen on that run. Loop the test a few thousand
times for confidence.
\end{warning}

\section{Coverage and confidence}

For a \maya{} app, aim for:

\begin{itemize}[leftmargin=*]
  \item 100\% statement coverage on \code{update}.
  \item Property tests for any invariant the model maintains.
  \item A handful of view snapshot tests for visual regressions.
  \item One or two end-to-end tests per major user flow.
\end{itemize}

The first two are cheap because \code{update} is pure. The last
two are the expensive ones; budget them deliberately.

\section{Pitfalls}

\begin{warning}
\textbf{Testing the renderer's exact output.} Different
terminals produce different bytes for the same content. Test
the content, not the bytes.
\end{warning}

\begin{warning}
\textbf{Snapshot fragility.} A snapshot that changes on every
unrelated UI tweak teaches the team to update snapshots without
reading them. Keep snapshots focused; cover one component at a
time.
\end{warning}

\begin{warning}
\textbf{Flaky timing tests.} If your end-to-end test sleeps for
real wall-clock time, it will be flaky on slow machines. Use
the virtual clock for any timing-dependent assertion.
\end{warning}

\begin{warning}
\textbf{Tests that depend on terminal capabilities.} A test
that checks ``truecolour bytes were emitted'' will fail on a
fake terminal with caps disabled. Either configure caps for
the test or assert against the more general behaviour.
\end{warning}

\section*{Workshop}

\begin{enumerate}[leftmargin=*]
  \item Write five unit tests for the two-counter app from
    Chapter~\ref{ch:elm}: each verifies one branch of
    \code{update}.
  \item Add a property test asserting that the counter never
    overflows under any sequence of messages.
  \item Configure your CI to fail on snapshot mismatches and to
    update snapshots only on a deliberate flag.
\end{enumerate}

\begin{recap}
The Elm architecture pays for itself most in tests.
\code{update} and \code{view} are pure and test directly with
function calls. The layout and serialise stages are pure too.
End-to-end tests use a fake terminal, an event queue, and a
virtual clock. Property tests cover the invariants the model
maintains. Concurrency tests need patience and a synchronous
pool.
\end{recap}
